\section{Introduction}
An upstream renewal desk can know an accepted obligation that a downstream execution gateway does not retain. When only a small service symbol crosses that boundary, the provider must choose both its terminal instruction alphabet and the branch-dependent encoder. This is a quantizer-design problem, but its distortion and feasible randomization are generated by a contract: an expected continuation promise, local participation limits, and the cost of providing intermediate rather than terminal service. The question is how much service value each additional retained symbol buys, and how that value changes when participation is checked before or after the symbol is drawn.

Finite-alphabet encoding, adjacent unbiased randomization, and matrix-search acceleration of ordered quantizer design are established ideas. In particular, \citet{Wu1991} already obtained $O(KN)$ work for optimal mean-square scalar quantization; stochastic rounding and dual quantization provide the relevant unbiased encoder primitives \citep{CrociEtAl2022,PagesWilbertz2012}. Our contribution is not to rename those primitives. It is to derive and solve the institution-generated quantization problem, including its costly service shortfall above the largest codeword, and to show which value comparisons survive when the aggregate promise is relaxed.

The analysis has three linked parts. First, for a saturated three-date renewal contract, deterministic execution gives an ordered one-sided quantizer, whereas expected participation permits adjacent unbiased lotteries. For heterogeneous quadratic costs, an optimal expected-participation book can be chosen with all nonhighest levels at branch caps and one continuously optimized highest level. The latter can lie strictly below the largest target and strictly between caps because intermediate service supplies the remaining obligation. A crossing-difference identity proves that the terminal array remains Monge after this continuous optimization. Classical favorable-staircase completion and SMAWK then yield $O(k)$ exact-rational oracle/comparison work per dynamic-programming layer. Pathwise participation eliminates the lottery advantage at the saturated endpoint for every alphabet budget.

Second, saturation does not confine the decision insight to a single root promise. We characterize the full-information service levels and minimum exact alphabet at every promise by a clipped common level. A same-alphabet scaling and saturation-lift argument bounds the change of every finite-budget value uniformly. Every strict expected-versus-pathwise advantage therefore persists on a certified interval of nonsaturated promises. An exact fixed-book reduction identifies the remaining joint promise-allocation problem rather than treating branch caps as targets after they cease to bind.

Third, we optimize bounded-overrun contracts over arbitrary finite certification catalogs with charges depending on the actual selected levels. An ordered-path dynamic program jointly chooses the codebook, service policy, and risk-feasible lotteries for every budget. This is more than pricing an already loss-minimizing book: an exact example changes the selected codewords when a single level becomes costly. The unrestricted continuous frontier and the charged finite-catalog frontier are separate, explicitly specified optimization problems.

The modeled renewal architecture is an operational abstraction, not an empirical claim about a named service provider. We use renewal language throughout and do not infer customer risk acceptance from the architecture. Exact regression, independently authored search substitution, reduced-network solver certificates, and computational measurements assess the algorithms; none substitutes for a proof of the contract reduction. A focused electronic companion contains the supporting proofs and implementation evidence. The complete preceding readers and broader compact-response theory remain independently accessible in the repository archive.
